\setchapterpreamble[u]{\margintoc}
\chapter{Super-resolution Microscopy}
\labch{SR-microscopy}

In the realm of optical microscopy, the pursuit of enhanced spatial resolution has led to the development of various 
    techniques collectively known as \ac{SR} microscopy.
The central aspect to the term \acl{SR} is the ability to surpass the classical established diffraction resolution
    limit by means of the technique.
\marginnote[*+8]{%
    An important fact that should be realized is that that mostly none of the \ac{SR} methods described here are useful in 
        other optical acquisition contexts then in microscopy.
    This is due to the property, that the methods rely on the ability to closely control the environment of the observation
        in various ways.
    As an example, \acs{SIM} is impossible to mimic, within the context of astronomy because we can not vary the
        illumination and in fact do not observe reflections of or emission due to illuminating the field.
}

\ac{SR} microscopy has been a term and a topic of study for almost half a century now.
In 1978, the first method to be developed, that may be categorized as a \ac{SR} method, is the so called
    \ac{4PIM} by \nameemph{Christoph} and \nameemph{Thomas Cremer} \sidecite{cremer1978b}.
It rethinks the classical design of sample illumination and attempts to limit the volume that is excited by the
    radiation within the sample space within one acquisition.
The light in the image space is thus reduced to the emission of the fluorescent molecules within this volume of the 
    sample and a reduction in uncertainty of the location within this volume of the sample is achieved.
By the repetition of these localized acquisitions, it is then possible to systematically cover the sample space and
    reconstruct the full sample estimate based on information from all of the acquisitions.
The realization that outperforming the classical diffraction limit by means of a specialized modality or setup of the 
    microscope is possible resulted in a rapid pace of innovation and lead to the invention of novel techniques, 
    altering the way of illuminating the sample and the preparation of the imaged sample itself.

A major breakthrough was presented in 1995 by \nameemph{J. Guerra}, who made use of a non-physical method to achieve 
    \ac{SR} in theory by including a modulation of the transfer function into the acquisition process and evaluating the 
    effect on the dependence of observed light in the image space on the source signal \sidecite{guerra1995}.
Later in 1999 \nameemph{Heintzmann} and \nameemph{Gustafsson} developed a deconvolution based reconstruction algorithm
    for \nameemph{Guerra's} method \sidecite{heintzmann1999}.
This change in perspective was also met with a large number of follow-ups which explored further approaches of
    utilizing the concept transfer function manipulation and pre-transfer inferred response signal control for the
    purpose of achieving \ac{SR}.

The culminated effort within the microscopy field lead to numerous advancements and unprecedented observations mostly in
    biology.
As a result and a appreciation was the granting of the Nobel Prize in Chemistry in 2014 shared three scientists who
    developed the \ac{STED} method \sidecite{weiss2014}.

\section{Comparison of \ac{SR} Methods}
\labsec{sr-comparison}

To this date, there are two major broad approaches of achieving \ac{SR} that have been demonstrated in practice 
    \sidecite{lal2016}.
In some modern cases a combination of methods one from each of the approach can be used alongside each other for
    achieving a greater resolution than one that is achievable by using them in isolation.
These broad categories are:
    \begin{description}
        \item[localization]%\labitm{localization}
            Concerns with reducing signal from the source structure to a limited volume or volumes that are sparse 
                within the image plane. 
            This bypasses the resolving limit using employing the aspect of the sensed signal, that there is assumed to 
                be little superposition from the individual fluorescent molecules.
            Lack of superposition can be used to distinguish individual emitting particles and if enough photons from
                one such particle is collected, its position can be determined with a low uncertainty.
        \item[modulation]%\labitm{modulation}
            By altering one of the aspect of generation of signal from the source structure, it induces a
                superposition of source signals caused by a range or discrete set of scales within the source structure.
            With a prior expectation or design of the set of scales superposed (aliased) and often by use of more than
                one such acquisition with different sets of scales admitted or included with a different factor within
                the superposition, the scales can be distinguished and mapped to the original scale of the source by
                reconstruction from the individual acquisitions.
    \end{description}

% FIX: Use references for the "localization" and "modulation" <kunzaatko> 
% https://tex.stackexchange.com/questions/1230/reference-name-of-description-list-item-in-latex
Usually a \emph{localization} based method uses illumination volume reduction for the purpose of limiting the volume of
    source structure originated signal in the sensed signal, it requires scanning or varying the location of the volume
    in a different manner in order to make a reconstruction of the full field source signal.
This is often at the demise of temporal resolution or alternatively an acquisition prioritizing temporal speed can
    severely reduce the resolution gain in the space axes.
In other words, fast localization \ac{SR} acquisitions usually do not retain their advertised \ac{SR} to the extent of
    still and long lasting acquisitions.
Aforementioned is a substantial case for using \emph{modulation} based methods for dynamic sources acquisitions.

The \emph{modulation} approach is a good fit for a dynamic (video) acquisition, since the photon efficiency relative to
    the efficiency of a standard acquisition is not lowered by a significant amount.
That is important due to the fact that if we want to acquire a live-cell for a period of time, it is desirable to limit
    the fluorophore bleaching caused by the cumulative effect of illumination during this time.
A popular choice for live-cell acquisitions is therefore \ac{SIM}.

A non-exhaustive table of \ac{SR} methods is presented in \reftab{sr-methods}.

% TODO: The table should not leave a half blank page. After the completed version of this chapter, the table should be
% moved to the place where a natural pagebreak would occur <08-01-24> 
\pagelayout{wide} % No margins
\begin{table*}
\vspace*{\fill}
    \RawFloats
    \centering
    \small
    \begin{tabularx}{\textwidth+\marginparsep+\marginparwidth}{l X c c c c}%{@{}l @{\extracolsep{\fill}} *{4}{l} @{}}
        \toprule
        \multirow[b]{2}{*}{\thead{\ac{SR}\\ Method}} & \multirow[b]{2}{*}{\thead{Method \\ Summary}}
                                                     & \multicolumn{3}{c}{\thead[b]{Resolution}}
                                                     & \multirow[b]{2}{*}{\thead{Year of \\ Origin}} \\
        \cmidrule(lr){3-5} 
                                                     & & {\centering\textbf{Lateral}} & {\centering\textbf{Axial}}
                                                     & {\centering\textbf{Temporal}} & \\
        \midrule
        \multirowcell{3}[0pt][l]{\acs{FIONA} \cite{yildiz2003}\\ \textbullet widefield \\ \textbullet localization } & {\footnotesize %
            Labels the sample with a single fluorophore within a diffraction radius. 
            Collects a sufficient number of photons from the isolated fluorophores to alleviate location noise and uses 
                curve fitting to find the center of the diffraction pattern that is assumed to be the true fluorophore 
            location.} & \(\qty{1.5}{\nano\meter}\) & \na & \(\qty{0.5}{\second}\) & 2003 \\
        \addlinespace
        \multirowcell{3}[0pt][l]{\acs{STED} \cite{hell1994} \\ \textbullet scanning \\ \textbullet localization} & {\footnotesize%
            Uses two laser beams for illumination: an excitation beam and a \acs{STED} beam. 
            The excitation beam excites the fluorophore photons and therefore causes fluorescence in the sample, while 
                the \acs{STED} beam, focused with small lateral offsets, depletes fluorescence in the outer regions of 
                the excitation spot reducing the volume where fluorophores are excited.%
            } & \(\qty{25}{\nano\meter}\) & \(\qty{100}{\nano\meter}\) & \(\qty{5}{\milli\second}\) & 1994 \\
        \addlinespace
        \multirowcell{3}[0pt][l]{\acs{PALM} \cite{betzig2006} \\ \textbullet widefield \\ \textbullet localization} & {\footnotesize%
            Stimulates the fluorophore labels in the sample with light and switching them from between the states of
                emitting and bleached, which is a state due to over activation of the fluorophore that wears off with 
                passing time.
            This is done in such a manner that the number of activated fluorophores is low and sparse within the sample.
            Similarly to \acs{FIONA}, this ensures their precise localization.
            It enough cycles are performed for a each set of fluorophores, a full image can be obtained by stacking the
                source estimations from each stack.
        } & \(\qty{2}{\nano\meter}\) & \(\qty{10}{\nano\meter}\) & \(\qty{5}{\second}\) & 2006 \\
        \addlinespace
        \multirowcell{3}[0pt][l]{\acs{STORM} \cite{rust2006,rahmanseresht2018} \\ \textbullet widefield \\ \textbullet  localization} & {\footnotesize%
                The idea is the same as in \acs{PALM}.
                Unlike \acs{PALM}, the sample is stained with specialized dyes that are photo-switchable.
                This means unlike \acs{PALM} no reliance is placed on the time it takes for the fluorophores to revive
                    from the bleached state and the switching may be done faster to the benefit of the temporal
                    resolution.
            } & \(\qty{4}{\nano\meter}\) & \(\qty{10}{\nano\meter}\) & \(\qty{1}{\second}\) & 2006  \\
        \addlinespace
        \multirowcell{3}[0pt][l]{\acs{4PIM} \cite{cremer1978b} \\ \textbullet scanning \\ \textbullet localization} & {\footnotesize%
            Uses the two opposing objective lenses focused on the same location within the sample.
            Further more, it uses a laser that is via reflection directed from both sides precisely aligned in relation
                to the sample and the reflected and emitted light can be collected leading to the ability of detecting
                the depth of the fluorophores by the effect of assessing, where relative to the laser beam wavelength
                the emission happens with better resolution.
                It improves only the axial resolution (using a classical wide-field microscope it is about
                    \(\qty{500}{\nano\meter}\)).
        } & \(\qty{180}{\nano\meter}\) & \(\qty{100}{\nano\meter}\) & \(\qty{50}{\milli\second}\) & 1978 \\
        \addlinespace
        \multirowcell{3}[0pt][l]{\acs{SIM} \cite{gustafsson2008, guerra1995, heintzmann1999} \\ \textbullet widefield \\ \textbullet modulation } & {\footnotesize%
            An illumination with non-uniform (modulated) intensity is used with sub-diffraction resolution lateral 
                variations.
            These variations modify the signal produced by the sample structure and shifts some of its variations that 
                are too rapid for the diffraction-limit and would not "survive" the optical system into lesser ones.
            This leads to aliasing within the sensed signal.
            It is possible to determine and the original locations of the variations, by an acquisition of multiple
                modulations and a reconstruction can be performed.
            % (Note, that variations are not meant in the literal sense of the mathematics definition, but rather
            %    variations as thought of intuitively)
        } & \(\qty{40}{\nano\meter}\) & \(\qty{300}{\nano\meter}\) & \(\qty{100}{\milli\second}\) & 1995 \\
        \bottomrule
    \end{tabularx}
    \normalsize
    \caption[skip=6pt]{%
        The most prolific\ac{SR} methods and their approximate optimal resolutions achievable in each axis as of January 
            2024.
        Keep in mind that the resolutions listed are optimal and are not, and often can not, be achieved in one 
            acquisition, because there is in most cases a compromise to be made. 
        If the resolving power for a single axis is optimized, the others suffer a loss in attainable resolution.
        An extreme example of this is \acs{PALM}, where the maximum frame rate is about \(\qty{12}{\per\minute}\) but an
            acquisition with the best lateral resolution of \(\qty{2}{\nano\meter}\) can take up to \(\qty{12}{\hour}\).
        \emph{Widefield} signifies that the acquisition of the full frame is performed in one take and \emph{scanning}
            that the acquisition of the full frame requires the accumulation of acquisition of multiple patches of
            different volumes of the sample space.
        Besides the impressive resolutions that are achieved by the \ac{SR} methods, it speaks volumes about their 
            inventors that they came up with some of the cheekiest acronyms widely used in contemporary science of 
            today.
    }
    \labtab{sr-methods}
\vspace*{\fill}
\end{table*}
\pagelayout{margin} % Restore margins

% TODO: Add labels <20-05-23> 
\section{Principle of SIM}
\labsec{sim-principle}

The principle which \ac{SIM} and other modulation based \ac{SR} methods stand on will be demonstrated by harmonic
    \ac{SIM}.
Harmonic \ac{SIM} alters the signal generated by the source structure via non-uniform illumination.
The illumination intensity takes the form of a \(2\)D harmonic function.

As a first step, let us consider, how a harmonic may affect aliasing of a signal in a one-dimensional example.
Let us have a model source structure, i.e. what we want to estimate, to be
    \begin{equation*}
        \source(x) = \cos(\omega x).
    \end{equation*}
Now, we want to imitate the illumination.
We can set the "illumination" as a harmonic in analog to the \(2\)D harmonic in \ac{SIM} that we want to get the 
    intuition for
    \begin{equation*}
        \ip(x) = \cos(\alpha\cdot\omega x),
    \end{equation*}
    where \(\alpha\) is a constant.
If we want to imitate \ac{SIM} \(\alpha \approx 1\) but \(\alpha \neq 1\).
To imitate the optical system, due to its linearity, it is adequate to perform the transfer of the "illumination", we
    multiply the source structure by it
    \begin{align*}
        \noisefree(x) &= \source(x) \cdot \ip(x) = \cos(\omega x)\cdot\cos(\alpha\omega x) \\[0.5em]
                      &= \cos\ab(\ab(\tfrac{\alpha + 1}{2} - \tfrac{\alpha - 1}{2})\cdot\omega
                      x)\cdot\cos\ab(\ab(\tfrac{\alpha + 1}{2} + \tfrac{\alpha - 1}{2})\cdot\omega x) \\[0.5em]
                      &= \frac{\cos\ab((1 + \alpha)\cdot\omega x) + \cos\ab((\alpha - 1) \cdot\omega x)}{2}
    \end{align*}
We can see from the above that some of the scales of variation of \(\source\)\sidenote[][*-20]{%
    Here, variations are not in the literal sense of the mathematics definition but rather intuitive variations as can
        be seen in \reffig{beats-mult}.
    Later in the \(2\)D case, we will do all the operations in Fourier space and the so called scales of variation will 
        be frequencies for us and the notion will stand on a firmer basis.
    } were approximately doubled (\(\alpha \approx 1\)) but some on the other hand are very close to \(0\).
This is the idea behind \ac{SIM}, where the variation that is lowered can now be observed.

% \pagebreak
\begin{marginfigure}[*-15]
    \marginpdftex{principle-of-sim-beats-mult}
    \caption{%
        We can see how the multiplication with the "illumination" affects the signal \(\noisefree\) that is generated. 
        Here, \(\alpha = 1.1\).
        It is clear from the graph, what we intuitively mean by the variation.
    }
    \labfig{beats-mult}
\end{marginfigure}

Even better than observing the exact case of a \(1\)D imitation of \ac{SIM} is to observe a physical analog in
    acoustics, that is often observed in real life.
If two harmonic acoustic waves with similar frequencies are superposed a \emph{beat} pattern emerges.
The real life use of this phenomenon is in sirens on an ambulance or on a police vehicle.
If we listen to them closely, it is possible to distinguish the high frequency harmonics and the low frequency beat
    sound.
A depiction of this phenomenon can be seen in \reffig{beats-add}.

\begin{marginfigure}[*-4]
    \marginpdftex{principle-of-sim-beats}
    \caption{Acoustic beat pattern that emerges when two close frequency harmonics are superposed.}
    \labfig{beats-add}
\end{marginfigure}

There is also an illustrative phenomenon that is commonplace in the real world of the above concept from \(1\)D in the 
    \(2\)D of \ac{SIM}.
It is in the form of \emph{moiré fringes} (or patterns).
They are present when two high-frequency patterns are superposed.
Examples of this in common use (\reffigshort{moire-real-life}) are:
    \begin{itemize}
        \item%
            A news anchorman with a striped shirt on an old low-resolution television.
            The patterns that are superposed are the stripes on the shirt and the scan-lines of the television.
        \item%
            A close-up photo of an \acs{LCD} or \acs{LED} screen  monitor. 
            The pattern are the pixels rows and columns of the camera and the monitor. 
            Scale of the patterns change when distance of the camera from the monitor.
        \item%
            The citizenship \acs{ID} card or a banknote, where moiré patterns are intensionally incorporated as
            a safety feature of the print.
    \end{itemize}

\begin{marginfigure}
    \centering
    \includegraphics[width=0.9\textwidth]{moire-art} \\[\medskipamount]
    
\includegraphics[width=0.9\textwidth]{moire-screen}
    \caption{%
        Moiré fringes are present when two high-frequency patterns are superposed.
        \vspace{-1.5em}
        \begin{center}
            {\scriptsize (courtesy of Emin Gabrielyan)}
        \end{center}
    }
    \labfig{moire-real-life}
\end{marginfigure}

\section{\ac{SIM}}%
\labsec{sim}

The aliasing in harmonic \ac{SIM} is better presented in the Fourier domain.
We will denote the Fourier images of functions by capital letters.
That is, we denote the Fourier image of the source \(\source\) as \(\F*{\source} = \Fsource\), the illumination
    \(\ip\) as \(\F*{\ip} = \Fip\) and so on.
This notation is also used in the experiment part.

The harmonic \ac{SIM} illumination pattern as already mentioned is the harmonic cosine function.
In \(2\)D, this is
% TODO: Add a figure of the illumination pattern here <08-01-24> 
    \begin{equation*}
        \labeq{ip-formula}
        \ip_{(m,\vec{k}_\ip,\phi)}(\vec{r}) = 1 + \tfrac{m}{2} \cos\ab(2\pi \vec{k}_\ip \cdot \vec{r} + \phi),
    \end{equation*}
    where \(m\), \(\vec{k}_\ip\) and \(\phi\) are constant parameters \emph{modulation}, \emph{phase-shift} (or
    \emph{orientation}) and \emph{phase-offset} of the illumination pattern respectively.
The Fourier image of a \(2\)D harmonic function is
    \newlength{\preeqlength}
    \settowidth{\preeqlength}{\(\displaystyle\mathsurround=0pt\F{\cos(2\pi \vec{k}_\ip\cdot \vec{r})}\)}
    \begin{equation*}
        \begin{split}
            \F{\cos(2\pi \vec{k}_\ip\cdot \vec{r})} &= \tfrac{1}{2}\F{\exp(2\pi \imag\vec{k}_\ip \cdot \vec{r}) + \exp(-2\pi
            \imag\vec{k}_\ip \cdot \vec{r})} \\[\medskipamount]
            &\begin{multlined}[t][\textwidth - \preeqlength]
            = \tfrac{1}{2} \iint_{\R^2} \ab(\exp(2\pi \imag\vec{k}_\ip \cdot \vec{r}) + \exp(-2\pi \imag\vec{k}_\ip \cdot
            \vec{r}))\\[-0.5em]
            \cdot\exp(-2\pi \imag \vec{k}\cdot \vec{r}) \odif{\vec{r}}
        \end{multlined}\\[\medskipamount]
            &\begin{multlined}[t][\textwidth - \preeqlength]
            =\tfrac{1}{2} \iint_{\R^2} \exp(2\pi \imag\vec{k}_\ip \cdot \vec{r}) \cdot \exp(-2\pi
                \imag \vec{k}\cdot \vec{r}) \odif{\vec{r}} \\
                + \tfrac{1}{2} \iint_{\R^2} \exp(-2\pi \imag\vec{k}_\ip \cdot \vec{r}) \cdot \exp(-2\pi \imag \vec{k}\cdot \vec{r}) \odif{\vec{r}}
            \end{multlined} \\[\medskipamount]
            &=\begin{multlined}[t][\textwidth - \preeqlength] 
                \frac{\F{\exp(2\pi \imag\vec{k}_\ip \cdot \vec{r})} + \F{\exp(-2\pi \imag\vec{k}_\ip \cdot \vec{r})}}{2}        
            \end{multlined} \\[\medskipamount]
            &= \tfrac{1}{2}\dirac\ab(\vec{k} + \vec{k}_\ip) + \tfrac{1}{2}\dirac\ab(\vec{k} - \vec{k}_{\ip}),
        \end{split}
    \end{equation*}
    which means that the Fourier image of the illumination pattern, due to linearity of \(\FourierSymbol\) is
    \begin{equation*}
      \F{\ip_{(m,\vec{k}_\ip,\phi)}(\vec{r})} = \Fip_{(m,\vec{k}_\ip,\phi)}(\vec{k}) = \frac{me^{-\imag \phi}}{4} \dirac_{-\vec{k}_\ip} + \dirac +
                          \frac{me^{\imag \phi}}{4} \dirac_{+\vec{k}_\ip}.
    \end{equation*}

We assume that the relation between excitation illumination intensity \(\ip\) and the emission intensity of a single 
    fluorophore is linear.
The relation between the emission intensity \(\noisefree\) at a point and the source structure \(\source\) (i.e. the 
    density of fluorescent particles) times the excitation intensity is also linear.
In fact we assume that the multiplicative coefficient is \(1\) without loss in generality since we are not interested in
    the actual number of photons produced, but instead the structure variations and so, in the spatial domain
    \begin{equation*}
      \noisefree(\vec{r}) = \source(\vec{r}) \cdot \ip(\vec{r}).
    \end{equation*}

We know that the in the Fourier domain, this translates to the convolution, i.e.
    \begin{equation*}
      \Fnoisefree(\vec{k}) = \Fsource(\vec{k}) * \Fip(\vec{k}).
    \end{equation*}

From this, it is clear that the structured illumination of the sample, induces spatial frequency aliasing of the sample
    source \(\source\) in the \(\noisefree\) signal.
This signal is then propagated through the optical system via the transfer function to produce the sensed signal
    \(\sensed\).
In the Fourier domain, this is expressed by
    \begin{equation*}
      \Fsensed(\vec{k}) = (\Fsource(\vec{k}) * \Fip(\vec{k})) \cdot \otf(\vec{k}).
    \end{equation*}
The reconstruction procedure then deals with antialiasing the frequencies of the sensed \(\Fsensed\) and placing them 
    into their correct locations within the reconstructed \(\Freconstructed\).
From \(\Freconstructed\), we can then obtain the reconstructed spatial source as 
    \begin{equation*}
        \Finverse*{\Freconstructed} = \reconstructed .
    \end{equation*}
This concept is further explained and performed in \refch{sim-reconstruction}. 

\section{Discussion}
\labsec{discussion}

Due to the fact that the only difference between the \ac{SIM} based acquisition and a regular acquisition is the
    non-uniform illumination, \ac{SIM} is composable with various setup modalities, available for regular view-field
    microscopy.
For example, it can be used in combination with \ac{TIRFM}, where noise can be reduced by reducing the illumination of
    the sample to a narrow axial slice.
The data that we are working with in the later chapters is acquired with this modality.

We are also not limited to the case of harmonic illumination.
There have been successful \ac{SIM} demonstrations with apriori unknown speckle patterns, with only the singular 
    condition, of summing into a uniform intensity \sidecite{yeh2017, mudry2012}.

It is also not limited to lateral resolution enhancement.
The harmonic patterns that we are using can also be expressed for axial slices that are not in one the focal plane and 
    the \(3\)D \ac{OTF} can be used in the reconstruction to enhance the image resolution in both axial and lateral 
    directions\sidecite{gustafsson2008}.
This approach requires \(15\) acquisitions with different parameters to in the reconstruction compared with \(9\) frame
    acquisitions of the standard Wiener filter based reconstruction for lateral resolution gain.

We are also not limited to a mere resolution doubling, as can be observed from \reftab{sr-methods} 
    \sidecite{gustafsson2005}.
It is essential though to realize  that the harmonic illumination pattern is limited in the intensity variation 
    frequency similarly by diffraction similarly to the wide way the optical system limits the resolution by
    diffraction\sidenote{%
        Thus \(\alpha \approx 1\) in \refsec{sim-principle} in not only the best option from the resolution gain 
        standpoint, but also a forced necessity in the \ac{SIM} realization.%
    }.
This problem is overcome by utilization of the non-linearity of the relation between excitation and emission near the
    saturation point of the fluorophores.
A non-linearity in the generation of photons by the fluorophores can be modelled and is desirable in this case, because
    it enables us to induce higher frequencies.
Let \(I_{\text{em}} = I_{\text{em}}(I_{\text{ex}})\) be the known non-linear relationship, where \(I_{\text{em}}\) is 
    the emission intensity and \(I_{\text{ex}}\) is the excitation intensity.
If we have the modulation \(m\) of \(I_{\text{ex}}\), such that the peak intensities of \(I_{\text{ex}}\) are above the
    saturation point of the fluorophores, we can express the relation \(I_{\text{em}}(I_{\text{ex}}(\vec{r}))\) by the
    sum of its harmonic components at a Fourier series.
The frequency components of the source \(\source\) corresponding to the Fourier series components of 
    \(I_{\text{em}}(I_{\text{ex}}(\vec{r}))\), including the higher frequencies than the ones included in the regular \ac{SIM}, are 
    aliased in the sensed image \(\Fsensed\). 
By multiple acquisitions and by using the calculated Fourier series coefficients, a chosen number of components can be 
    linearly separated similarly as in \refsec{sim-reconstruction}, and the reconstruction \(\Freconstructed\) can be
    assembled from these components.
If you think about the aforementioned procedure, the lateral resolution that can be achieved is theoretically unlimited.
What is the problem\sidenote{%
        There must be one, because the prospect of unlimited resolution is unthinkable.
    }?
If we were to expand the relation \(I_{\text{em}}(I_{\text{ex}}(\vec{r}))\) into its Fourier series, we would observe that the higher
    of the harmonic component \(\zeta_k(\vec{r})\) the lower the corresponding coefficient \(\rho_k\) and at some point, the 
    component would become undistinguishable from the photon noise.
This approach has also been implemented in combination with the photo-switchable dyes from \acs{STORM}, which achieved
    the lateral resolution of \(\qty{50}{\nano\meter}\) \sidecite{rego2012}.

% TODO: prakash2021 is good for an overview and the discussion of ways forward <07-01-24> 
